% ==================================================
% Supplement: Local-global + arithmetic refinements for the
% universal anomaly comparison. Five refinement cycles inscribed:
% (1) Atiyah scope separation; (2) Koszul locus chart vs moduli;
% (3) Two loci intersection via Kuga--Satake period map;
% (4) Nine Heegner d_K as conditional rank-8 arithmetic problem;
% (5) Motivic Galois scope of Bridgeland finiteness.
% ==================================================

\section{Local-global scope and arithmetic refinements}
\label{sec:local-global-arithmetic}

Five refinements sharpen the scope of the conditional bridge. Each names a chart-local/global dichotomy or an arithmetic-theoretic refinement that the main statements assume implicitly.

\subsection{Atiyah scope separation: four global statements that admit no chart-local avatar}
\label{subsec:atiyah-scope-separation}

The four statements
\begin{itemize}
  \item the Hodge location of Corollary~\ref{cor:hodge}, namely $[m_3, B^{(2)}]_X \in H^2(X,\Omega^1_X)=H^{1,2}(X)$;
  \item the Yukawa non-vanishing $\int_{Q_5} \at_{Q_5}^{\cup 3} \neq 0$ of Proposition~\ref{prop:quintic-nonvanish};
  \item the K\"unneth additivity $\at_{K3 \times E} = \at_{K3} \boxplus \at_E$ of Proposition~\ref{prop:k3e-vanish};
  \item the vanishing $\alpha_X=0 \in H^2(X, \Omega^1_X)$ defining the Atiyah--Connes zero-locus candidate of Definition~\ref{def:atiyah-connes-zero-locus}
\end{itemize}
each live strictly at \emph{global} scope on $X$. On any single affine chart $U_\alpha \subset X$ trivialising $T_X$ --- witnessed by a chart-local Chern connection $\nabla_\alpha$ --- the restriction $\at_\cE|_{U_\alpha} \in \Ext^1_{U_\alpha}(\cE|_{U_\alpha}, \cE|_{U_\alpha} \otimes \Omega^1_{U_\alpha})$ vanishes identically: $\nabla_\alpha$ is a primitive, and every chart-local formula written in local coordinates with a named $\nabla_\alpha$ is tautologous at the level of $\at$. The four global statements are not shadowed by chart-local computations; they are born from patching combinatorics.

Conversely, the chain-level constructions of §\ref{sec:cocycle} and Theorem~\ref{thm:clemens-extension} are intrinsically chart-local. The Massey primitive $\alpha$ with $d\alpha = f \circ g$ exists chart-locally; the Kotake--Kato product heat-kernel is chart-local; the Costello--Li parametrix pinning is chart-local on each Kuranishi chart. The content of each vanishing theorem is the obstruction to patching these chart-local primitives into one global primitive. Local vanishing is automatic. Global vanishing is the theorem.

\subsection{Koszul locus: chart-local on $X$ versus global in $\cM_{\mathrm{cx}}(X)$}
\label{subsec:koszul-scope}

The locus $\cU^{\mathrm{adm}}(X) \subset \cM_{\mathrm{cx}}(X)$ is a subscheme of the Kuranishi moduli space of complex structures on the underlying smooth manifold, not a subset of $X$. Two equivalent global definitions serve. Globally on $\cM_{\mathrm{cx}}(X)$, the relative cocycle
\[
  [m_3, B^{(2)}]_{X/\cM_{\mathrm{cx}}(X)} \;\in\; R^2 (\pi_\cM)_* \Omega^1_{X_{\mathrm{univ}}/\cM_{\mathrm{cx}}(X)}
\]
on the universal family $\pi_\cM : X_{\mathrm{univ}} \to \cM_{\mathrm{cx}}(X)$ defines a coherent sheaf. The vanishing locus of the corresponding section is the obstruction-theoretic candidate for $\cU^{\mathrm{adm}}(X)$; any openness statement requires a separate smoothness or constant-rank hypothesis. Equivalently, $\cU^{\mathrm{adm}}(X)$ parametrises $[J] \in \cM_{\mathrm{cx}}(X)$ for which the Positselski contracting homotopy $h$ with $[d, h] = \id - p$ exists at every order of the completed modular cyclic obstruction filtration. Theorem~\ref{thm:main} makes the equivalence of these two definitions conditional on the bridge map \textup{(B1)}.

A Kuranishi chart $V_\beta \subset \cM_{\mathrm{cx}}(X)$ is \emph{not} a chart $U_\alpha \subset X$: the former is an open subset of moduli space parametrising a Kodaira--Spencer slice at some $[J_\beta]$; the latter is an open subset of the threefold. Two propagations must not be conflated: chart-overlap Čech patching on $U_\alpha \cap U_{\alpha'} \subset X$ of §\ref{subsec:atiyah-local-global} (gluing local connections into $\at_X$), and chart-overlap gluing on $V_\beta \cap V_{\beta'} \subset \cM_{\mathrm{cx}}$ (gluing Kuranishi obstruction data into $\cU^{\mathrm{adm}}$). Openness of $\cU^{\mathrm{adm}}$ is the Kodaira--Spencer compatibility of chart-uniform Koszulness witnesses across moduli-chart overlaps.

The three witnesses of §\ref{sec:clemens} live on three distinct moduli components: $\cM_{\mathrm{cx}}(Q_5)$ is $101$-dimensional and is a falsifier for the unconditional pentagon (Corollary~\ref{cor:quintic-fails}); $\cM_{\mathrm{cx}}(K3 \times E)$ is $21$-dimensional, but its admissible locus is conditional on the scalar computation $\lambda_{\mathrm{Cyc}}=0$ and on any Borcea--Voisin recomputation (Corollary~\ref{cor:k3e-locus}); the Clemens stratum lies in a non-K\"ahler component on which Dolbeault cohomology is ill-defined, and $\cU^{\mathrm{adm}}$ is given by the Fr\"olicher-degenerate reformulation of Theorem~\ref{thm:clemens-extension}.

\subsection{Two Koszul loci intersection: Atiyah side versus Humbert side via Kuga--Satake}
\label{subsec:two-loci-kuga-satake}

Theorem~\ref{thm:two-loci} intersects two a-priori-unrelated Koszul loci. The \emph{Atiyah-cocycle locus} $\cU^{\mathrm{adm}}_{\at}(X) \subset \cM_{\mathrm{cx}}(X)$ is the moduli-global vanishing locus of $[m_3, B^{(2)}]_{X/\cM_{\mathrm{cx}}(X)}$ of §\ref{subsec:koszul-scope}, an Atiyah-theoretic condition on $X$ itself. The \emph{Humbert-period locus} $\cP^{-1}(U^{\mathrm{adm}}) \subset \cM_{\mathrm{cx}}(X)$ is the preimage of the Kuga--Satake period image of the Bruinier admissible locus, a period-theoretic condition on $\overline{\cM_{\mathrm{A_2}}} = \Gamma \backslash D$ (the arithmetic quotient of the period domain, a Shimura variety).

The two a-priori-unrelated loci intersect in the stage-one canonical locus: this is the content of Theorem~\ref{thm:two-loci}. The period map $\cP$ reads Hodge-theoretic data; the Atiyah cocycle reads local Kuranishi-obstruction data. Their intersection equals $\cU^{\mathrm{canonical}}_{\Phi_3}$ by stage-one canonicality plus Griffiths transversality.

The Kuga--Satake period map (Kuga--Satake 1967, Morrison 1984 Theorem~7.9) is the embedding $\mathrm{KS} : D \to \cA_g^{\mathrm{KS}}$ lifting the K3 period domain $D = \mathrm{O}^+(2, 20)/(\mathrm{SO}(2) \times \mathrm{O}(20))$ to the Siegel modular variety $\cA_g^{\mathrm{KS}}$ of principally polarised abelian varieties of dimension $g = 2^{19}$. The lift is Hodge-compatible: the $(2,0)$-part of a K3 becomes the holomorphic $1$-forms of its Kuga--Satake abelian variety. Humbert divisors on $\cA_g^{\mathrm{KS}}$ classify abelian varieties whose endomorphism ring contains a real quadratic order; their preimages under $\mathrm{KS}$ and $\cP$ form the loci on which Bruinier's positivity cut imposes a definite signature at the Mukai-enhanced central charge $K = 8$. The passage through Kuga--Satake is necessary: the same positivity condition on the K3 side involves the transcendental lattice, which is directly readable on the abelian variety side via endomorphism structure. The factorisation $\cP = \mathrm{KS} \circ \cP_{K3}$ for K3-fibred $X$ identifies $\cU^{\mathrm{canonical}}_{\Phi_3}(X)$ as the Atiyah condition \emph{intersected with} the period image of the Bruinier locus: two independent conditions, the second Kuga--Satake-arithmetic.

\subsection{Nine Heegner discriminants: conditional rank-$8$ arithmetic problem}
\label{subsec:heegner-chenevier}

Proposition~\ref{prop:heegner} retains
$d_K \in \{3, 4, 7, 8, 11, 15, 19, 20, 24\}$ as a conditional
rank-$8$ arithmetic candidate. As currently justified, this set is the
initial segment $d\leq24$ of fundamental negative discriminants with
class number at most $2$, not a proved Bruinier--Chenevier selection.
The classical class-number-one fundamental discriminants remain
\[
  \{3,4,7,8,11,19,43,67,163\}.
\]

Chenevier's determinants and Bruinier's Heegner-divisor Chern-class
formula are relevant technologies, but they do not by themselves force
$K=8$, $\hbar^2=-1/8$, or the nine-element list. To make the refinement
theorem-grade one must define an explicit rank-$8$ admissibility
predicate $\mathcal C_8(d)$, compute the Bruinier coefficient/sign in
that predicate, and prove that its solution set is exactly the stated
nine discriminants.

The nine discriminants $\{3, 4, 7, 8, 11, 15, 19, 20, 24\}$ should not
be justified by the unit-group condition: every imaginary quadratic
order contains $\{\pm1\}$, while only $d_K=3,4$ have exceptional extra
units. Nor should Chenevier determinants be cited as selecting the list
without a residual pseudodeformation datum and a comparison map.

Thus the arithmetic refinement is conditional: either the
rank-$8$ predicate and Bruinier coefficient calculation must be
supplied, or the statement should revert to the classical Heegner list
with a separate positivity problem.

\subsection{Motivic Galois scope of Bridgeland finiteness}
\label{subsec:motivic-galois-scope}

Corollary~\ref{cor:bridgeland}'s conditional Bridgeland finiteness sits inside a larger arithmetic scope governed by the motivic Galois group. The \emph{motivic Galois group} $G_{\mathrm{mot}}(X)$ of a Calabi--Yau threefold $X$ is the Tannakian dual of the category of mixed motives over $X$ (Deligne--Milne 1982, Andr\'e 1996). Conjecturally, $G_{\mathrm{mot}}(X)$ is a reductive linear algebraic group acting on $\bigoplus_n H^n_{\mathrm{mot}}(X) \cong \bigoplus_n H^n_{\mathrm{dR}}(X)$ by restrictions of the Mumford--Tate group of the complex Hodge structure.

The Mumford--Tate--Galois--Sato--Lannes principle, in this manuscript's
working nomenclature, asserts that for a Calabi--Yau threefold
admitting a K3 fibration, the motivic Galois group of the Kuga--Satake
lift factors through the orthogonal group of the Mukai lattice:
$G_{\mathrm{mot}}(\mathrm{KS}(X)) \subseteq \mathrm{O}(L_{K3,\Qbb})$.
This is a programme-level conjectural input, not a
Chenevier--Lannes theorem.

The Bridgeland finiteness statement of Corollary~\ref{cor:bridgeland} is conditional on Bridgeland's conjecture that $\Stab(D^b(X))$ is a finite-dimensional complex manifold with locally-closed period maps. The arithmetic refinement: if MGSL holds for $X$, then the Koszul-admissible locus $\cU^{\mathrm{adm}}(X) \cap \cM_{\mathrm{cx}}(X)$ decomposes into $G_{\mathrm{mot}}(X)$-orbits under the motivic Galois action on the period domain; each orbit is a Shimura subvariety, and Bridgeland finiteness on each orbit follows from the arithmetic-volume finiteness of Shimura subvarieties (Zhang 2014). The motivic-Galois conditional statement is therefore strictly stronger than the Bridgeland conditional: MGSL $\Rightarrow$ Bridgeland on $\cU^{\mathrm{adm}}(X)$ but not conversely.

Two scope-separations are load-bearing. First, the motivic Galois scope applies only to $X$ in the K3-fibered regime where Kuga--Satake lifts exist; Clemens' non-K\"ahler stratum has no motivic Galois group in the conjectural sense (the mixed motives of a non-K\"ahler threefold lack the Hodge structure required for the Tannakian construction). Second, the Bridgeland scope applies more broadly but is weaker: it gives finiteness on compact subsets of $\Stab(D^b(X))$ but not global decomposition into Shimura subvarieties. Theorem~\ref{thm:main} holds in both scopes; its refinements via Heegner discriminants (Proposition~\ref{prop:heegner}) and two-loci intersection (Theorem~\ref{thm:two-loci}) use Kuga--Satake and hence live in the motivic-Galois scope.

\begin{remark}[Bibliographic anchors]
\label{rem:arithmetic-biblio}
The arithmetic refinements invoke: Chenevier 2014 (\emph{Astérisque} 324, determinants of Galois representations), Deligne--Milne 1982 (Tannakian categories), Andr\'e 1996 (motivic Galois group), Chenevier--Lannes 2019 (mass formula for Niemeier lattices with $M_{23}$ twist), Morrison 1984 Theorem 7.9 (Kuga--Satake construction), Morrison 1984 Theorem 8.1 ($K3 \times E$ motivic Galois), Zhang 2014 (arithmetic volume of Shimura subvarieties). None enter the proof of Theorem~\ref{thm:main}; they refine the scope of the corollaries (Theorem~\ref{thm:two-loci}, Proposition~\ref{prop:heegner}, Corollary~\ref{cor:bridgeland}, Corollary~\ref{cor:m23}) into arithmetic-precise statements.
\end{remark}
